\section{Related Works}
\label{sec:related_works}

\paragraph{Reinforcement learning for LLM post-training}
RL has become a prominent paradigm for LLM post-training, especially for improving reasoning, coding, and long-horizon tasking capabilities.
A major recent direction is reinforcement learning with verifiable rewards (RLVR), where outcome rewards can be obtained from automatically checkable signals such as mathematical answers, code execution, or task success.
Representative works include GRPO~\citep{Shao24GRPO}, DAPO~\citep{Yu2025DAPO}, and GSPO~\citep{gspo}, which develop scalable RL recipes for mathematical and reasoning tasks.
In parallel, recent studies have also explored off-policy or data-reuse variants of reinforced fine-tuning~\citep{remix,zhan2026exgrpo}, highlighting the importance of correcting distributional mismatch during LLM RL.
More recently, RL has also been adopted in agentic model training, where models learn from environment feedback in tool-use, coding, and multi-turn interaction settings.
Examples include GLM-4.5~\citep{glm}, Kimi K2.5~\citep{kimiteam2026kimik25visualagentic}, and ROME~\citep{ROME}.
Together, these works show the broad potential of RL for improving LLM capabilities, while also highlighting the fragility of RL training under long-horizon rollouts, complex optimization dynamics, and system-level implementation details.

\vspace{-0.2cm}
\paragraph{Training-inference mismatch in LLM RL}
Recent work has identified training-inference mismatch as an important source of instability in modern LLM RL systems.
A common line of algorithmic approaches incorporates sampler-side information into the training update.
For example, TIS~\citep{mismatch1} uses the training-to-sampler probability ratio as a clipped correction weight, while MIS~\citep{mismatch2} filters tokens or sequences with extreme mismatch signals.
Learning rate decay~\citep{lrdecay} provides another optimization-side remedy by reducing the update magnitude when mismatch-related instability emerges.
In parallel, infrastructure-level approaches study how system design choices affect training-inference mismatch. 
For example, \citet{fp16} show that numerical precision can be a direct source of mismatch and propose using FP16 rollout to reduce the discrepancy between training and inference engines. 
\citet{li2026qurl} further investigate low-precision RL systems and analyze how quantized rollout affects reasoning-oriented RL training. 
These works reduce mismatch from the system side, while sampler-aware methods such as TIS and MIS correct or filter the training update from the algorithm side. 
Our work is complementary: instead of only reducing mismatch, we study its objective-level consequence and ask whether a synchronized update should be accepted as the next inference policy.